\documentclass[12pt]{article}

% Idioma y codificación
\usepackage[spanish, es-tabla]{babel}       %es-tabla para que se titule "Tabla"
\usepackage[utf8]{inputenc}

% Márgenes
\usepackage[a4paper,top=3cm,bottom=2.5cm,left=3cm,right=3cm]{geometry}

% Comentarios de bloque
\usepackage{verbatim}

% Paquetes de links
\usepackage[hidelinks]{hyperref}    % Permite enlaces
\usepackage{url}                    % redirecciona a la web

% Más opciones para enumeraciones
\usepackage{enumitem}

% Personalizar la portada
\usepackage{titling}


% Paquetes de tablas
\usepackage{multirow}


%------------------------------------------------------------------------

%Paquetes de figuras
\usepackage{caption}
\usepackage{subcaption} % Figuras al lado de otras
\usepackage{float}      % Poner figuras en el sitio indicado H.


% Paquetes de imágenes
\usepackage{graphicx}       % Paquete para añadir imágenes
\usepackage{transparent}    % Para manejar la opacidad de las figuras

% Paquete para usar colores
\usepackage[dvipsnames]{xcolor}
\usepackage{pagecolor}      % Para cambiar el color de la página

% Habilita tamaños de fuente mayores
\usepackage{fix-cm}


%------------------------------------------------------------------------

% Paquetes de matemáticas
\usepackage{mathtools, amsfonts, amssymb, mathrsfs}
\usepackage[makeroom]{cancel}     % Simplificar tachando
\usepackage{polynom}    % Divisiones y Ruffini
\usepackage{units} % Para poner fracciones diagonales con \nicefrac

\usepackage{pgfplots}   %Representar funciones
\pgfplotsset{compat=1.18}  % Versión 1.18

\usepackage{tikz-cd}    % Para usar diagramas de composiciones
\usetikzlibrary{calc}   % Para usar cálculo de coordenadas en tikz

%Definición de teoremas, etc.
\usepackage{amsthm}
%\swapnumbers   % Intercambia la posición del texto y de la numeración

\theoremstyle{plain}

\makeatletter
\@ifclassloaded{article}{
  \newtheorem{teo}{Teorema}[section]
}{
  \newtheorem{teo}{Teorema}[chapter]  % Se resetea en cada chapter
}
\makeatother

\newtheorem{coro}{Corolario}[teo]           % Se resetea en cada teorema
\newtheorem{prop}[teo]{Proposición}         % Usa el mismo contador que teorema
\newtheorem{lema}[teo]{Lema}                % Usa el mismo contador que teorema

\theoremstyle{remark}
\newtheorem*{observacion}{Observación}

\theoremstyle{definition}

\makeatletter
\@ifclassloaded{article}{
  \newtheorem{definicion}{Definición} [section]     % Se resetea en cada chapter
}{
  \newtheorem{definicion}{Definición} [chapter]     % Se resetea en cada chapter
}
\makeatother

\newtheorem*{notacion}{Notación}
\newtheorem*{ejemplo}{Ejemplo}
\newtheorem*{ejercicio*}{Ejercicio}             % No numerado
\newtheorem{ejercicio}{Ejercicio} [section]     % Se resetea en cada section


% Modificar el formato de la numeración del teorema "ejercicio"
\renewcommand{\theejercicio}{%
  \ifnum\value{section}=0 % Si no se ha iniciado ninguna sección
    \arabic{ejercicio}% Solo mostrar el número de ejercicio
  \else
    \thesection.\arabic{ejercicio}% Mostrar número de sección y número de ejercicio
  \fi
}


% \renewcommand\qedsymbol{$\blacksquare$}         % Cambiar símbolo QED
%------------------------------------------------------------------------

% Paquetes para encabezados
\usepackage{fancyhdr}
\pagestyle{fancy}
\fancyhf{}

\newcommand{\helv}{ % Modificación tamaño de letra
\fontfamily{}\fontsize{12}{12}\selectfont}
\setlength{\headheight}{15pt} % Amplía el tamaño del índice


%\usepackage{lastpage}   % Referenciar última pag   \pageref{LastPage}
\fancyfoot[C]{\thepage}

%------------------------------------------------------------------------

% Conseguir que no ponga "Capítulo 1". Sino solo "1."
\makeatletter
\@ifclassloaded{book}{
  \renewcommand{\chaptermark}[1]{\markboth{\thechapter.\ #1}{}} % En el encabezado
    
  \renewcommand{\@makechapterhead}[1]{%
  \vspace*{50\p@}%
  {\parindent \z@ \raggedright \normalfont
    \ifnum \c@secnumdepth >\m@ne
      \huge\bfseries \thechapter.\hspace{1em}\ignorespaces
    \fi
    \interlinepenalty\@M
    \Huge \bfseries #1\par\nobreak
    \vskip 40\p@
  }}
}
\makeatother

%------------------------------------------------------------------------
% Paquetes de cógido
\usepackage{minted}
\renewcommand\listingscaption{Código fuente}

\usepackage{fancyvrb}
% Personaliza el tamaño de los números de línea
\renewcommand{\theFancyVerbLine}{\small\arabic{FancyVerbLine}}

% Estilo para C++
\newminted{cpp}{
    frame=lines,
    framesep=2mm,
    baselinestretch=1.2,
    linenos,
    escapeinside=||
}



\usepackage{listings} % Para incluir código desde un archivo

\renewcommand\lstlistingname{Código Fuente}
\renewcommand\lstlistlistingname{Índice de Códigos Fuente}

% Definir colores
\definecolor{vscodepurple}{rgb}{0.5,0,0.5}
\definecolor{vscodeblue}{rgb}{0,0,0.8}
\definecolor{vscodegreen}{rgb}{0,0.5,0}
\definecolor{vscodegray}{rgb}{0.5,0.5,0.5}
\definecolor{vscodebackground}{rgb}{0.97,0.97,0.97}
\definecolor{vscodelightgray}{rgb}{0.9,0.9,0.9}

% Configuración para el estilo de C similar a VSCode
\lstdefinestyle{vscode_C}{
  backgroundcolor=\color{vscodebackground},
  commentstyle=\color{vscodegreen},
  keywordstyle=\color{vscodeblue},
  numberstyle=\tiny\color{vscodegray},
  stringstyle=\color{vscodepurple},
  basicstyle=\scriptsize\ttfamily,
  breakatwhitespace=false,
  breaklines=true,
  captionpos=b,
  keepspaces=true,
  numbers=left,
  numbersep=5pt,
  showspaces=false,
  showstringspaces=false,
  showtabs=false,
  tabsize=2,
  frame=tb,
  framerule=0pt,
  aboveskip=10pt,
  belowskip=10pt,
  xleftmargin=10pt,
  xrightmargin=10pt,
  framexleftmargin=10pt,
  framexrightmargin=10pt,
  framesep=0pt,
  rulecolor=\color{vscodelightgray},
  backgroundcolor=\color{vscodebackground},
}

%------------------------------------------------------------------------

% Comandos definidos
\newcommand{\bb}[1]{\mathbb{#1}}
\newcommand{\cc}[1]{\mathcal{#1}}

% I prefer the slanted \leq
\let\oldleq\leq % save them in case they're every wanted
\let\oldgeq\geq
\renewcommand{\leq}{\leqslant}
\renewcommand{\geq}{\geqslant}

% Si y solo si
\newcommand{\sii}{\iff}

% Letras griegas
\newcommand{\eps}{\epsilon}
\newcommand{\veps}{\varepsilon}
\newcommand{\lm}{\lambda}

\newcommand{\ol}{\overline}
\newcommand{\ul}{\underline}
\newcommand{\wt}{\widetilde}
\newcommand{\wh}{\widehat}

\let\oldvec\vec
\renewcommand{\vec}{\overrightarrow}

% Derivadas parciales
\newcommand{\del}[2]{\frac{\partial #1}{\partial #2}}
\newcommand{\Del}[3]{\frac{\partial^{#1} #2}{\partial^{#1} #3}}
\newcommand{\deld}[2]{\dfrac{\partial #1}{\partial #2}}
\newcommand{\Deld}[3]{\dfrac{\partial^{#1} #2}{\partial^{#1} #3}}


\newcommand{\AstIg}{\stackrel{(\ast)}{=}}
\newcommand{\Hop}{\stackrel{L'H\hat{o}pital}{=}}

\newcommand{\red}[1]{{\color{red}#1}} % Para integrales, destacar los cambios.

% Método de integración
\newcommand{\MetInt}[2]{
    \left[\begin{array}{c}
        #1 \\ #2
    \end{array}\right]
}

% Declarar aplicaciones
% 1. Nombre aplicación
% 2. Dominio
% 3. Codominio
% 4. Variable
% 5. Imagen de la variable
\newcommand{\Func}[5]{
    \begin{equation*}
        \begin{array}{rrll}
            #1:& #2 & \longrightarrow & #3\\
               & #4 & \longmapsto & #5
        \end{array}
    \end{equation*}
}

%------------------------------------------------------------------------
\usepgfplotslibrary{fillbetween}
\DeclareMathOperator{\Var}{Var}
\DeclareMathOperator{\Cov}{Cov}

\begin{document}

    % 1. Foto de fondo
    % 2. Título
    % 3. Encabezado Izquierdo
    % 4. Color de fondo
    % 5. Coord x del titulo
    % 6. Coord y del titulo
    % 7. Fecha

    
    % 1. Foto de fondo
% 2. Título
% 3. Encabezado Izquierdo
% 4. Color de fondo
% 5. Coord x del titulo
% 6. Coord y del titulo
% 7. Fecha

\newcommand{\portada}[7]{

    \portadaBase{#1}{#2}{#3}{#4}{#5}{#6}{#7}
    \portadaBook{#1}{#2}{#3}{#4}{#5}{#6}{#7}
}

\newcommand{\portadaExamen}[7]{

    \portadaBase{#1}{#2}{#3}{#4}{#5}{#6}{#7}
    \portadaArticle{#1}{#2}{#3}{#4}{#5}{#6}{#7}
}




\newcommand{\portadaBase}[7]{

    % Tiene la portada principal y la licencia Creative Commons
    
    % 1. Foto de fondo
    % 2. Título
    % 3. Encabezado Izquierdo
    % 4. Color de fondo
    % 5. Coord x del titulo
    % 6. Coord y del titulo
    % 7. Fecha
    
    
    \thispagestyle{empty}               % Sin encabezado ni pie de página
    \newgeometry{margin=0cm}        % Márgenes nulos para la primera página
    
    
    % Encabezado
    \fancyhead[L]{\helv #3}
    \fancyhead[R]{\helv \nouppercase{\leftmark}}
    
    
    \pagecolor{#4}        % Color de fondo para la portada
    
    \begin{figure}[p]
        \centering
        \transparent{0.3}           % Opacidad del 30% para la imagen
        
        \includegraphics[width=\paperwidth, keepaspectratio]{assets/#1}
    
        \begin{tikzpicture}[remember picture, overlay]
            \node[anchor=north west, text=white, opacity=1, font=\fontsize{60}{90}\selectfont\bfseries\sffamily, align=left] at (#5, #6) {#2};
            
            \node[anchor=south east, text=white, opacity=1, font=\fontsize{12}{18}\selectfont\sffamily, align=right] at (9.7, 3) {\textbf{\href{https://losdeldgiim.github.io/}{Los Del DGIIM}}};
            
            \node[anchor=south east, text=white, opacity=1, font=\fontsize{12}{15}\selectfont\sffamily, align=right] at (9.7, 1.8) {Doble Grado en Ingeniería Informática y Matemáticas\\Universidad de Granada};
        \end{tikzpicture}
    \end{figure}
    
    
    \restoregeometry        % Restaurar márgenes normales para las páginas subsiguientes
    \pagecolor{white}       % Restaurar el color de página
    
    
    \newpage
    \thispagestyle{empty}               % Sin encabezado ni pie de página
    \begin{tikzpicture}[remember picture, overlay]
        \node[anchor=south west, inner sep=3cm] at (current page.south west) {
            \begin{minipage}{0.5\paperwidth}
                \href{https://creativecommons.org/licenses/by-nc-nd/4.0/}{
                    \includegraphics[height=2cm]{assets/Licencia.png}
                }\vspace{1cm}\\
                Esta obra está bajo una
                \href{https://creativecommons.org/licenses/by-nc-nd/4.0/}{
                    Licencia Creative Commons Atribución-NoComercial-SinDerivadas 4.0 Internacional (CC BY-NC-ND 4.0).
                }\\
    
                Eres libre de compartir y redistribuir el contenido de esta obra en cualquier medio o formato, siempre y cuando des el crédito adecuado a los autores originales y no persigas fines comerciales. 
            \end{minipage}
        };
    \end{tikzpicture}
    
    
    
    % 1. Foto de fondo
    % 2. Título
    % 3. Encabezado Izquierdo
    % 4. Color de fondo
    % 5. Coord x del titulo
    % 6. Coord y del titulo
    % 7. Fecha


}


\newcommand{\portadaBook}[7]{

    % 1. Foto de fondo
    % 2. Título
    % 3. Encabezado Izquierdo
    % 4. Color de fondo
    % 5. Coord x del titulo
    % 6. Coord y del titulo
    % 7. Fecha

    % Personaliza el formato del título
    \pretitle{\begin{center}\bfseries\fontsize{42}{56}\selectfont}
    \posttitle{\par\end{center}\vspace{2em}}
    
    % Personaliza el formato del autor
    \preauthor{\begin{center}\Large}
    \postauthor{\par\end{center}\vfill}
    
    % Personaliza el formato de la fecha
    \predate{\begin{center}\huge}
    \postdate{\par\end{center}\vspace{2em}}
    
    \title{#2}
    \author{\href{https://losdeldgiim.github.io/}{Los Del DGIIM}}
    \date{Granada, #7}
    \maketitle
    
    \tableofcontents
}




\newcommand{\portadaArticle}[7]{

    % 1. Foto de fondo
    % 2. Título
    % 3. Encabezado Izquierdo
    % 4. Color de fondo
    % 5. Coord x del titulo
    % 6. Coord y del titulo
    % 7. Fecha

    % Personaliza el formato del título
    \pretitle{\begin{center}\bfseries\fontsize{42}{56}\selectfont}
    \posttitle{\par\end{center}\vspace{2em}}
    
    % Personaliza el formato del autor
    \preauthor{\begin{center}\Large}
    \postauthor{\par\end{center}\vspace{3em}}
    
    % Personaliza el formato de la fecha
    \predate{\begin{center}\huge}
    \postdate{\par\end{center}\vspace{5em}}
    
    \title{#2}
    \author{\href{https://losdeldgiim.github.io/}{Los Del DGIIM}}
    \date{Granada, #7}
    \thispagestyle{empty}               % Sin encabezado ni pie de página
    \maketitle
    \vfill
}
    \portadaExamen{ffccA4.jpg}{Probabilidad\\Examen VI}{Probabilidad. Examen VI}{MidnightBlue}{-8}{28}{2024-2025}{Arturo Olivares Martos \\ José Juan Urrutia Milán}

    \begin{description}
        \item[Asignatura] Probabilidad.
        \item[Curso Académico] 2019-20.
        \item[Grado] Grado en Matemáticas y Doble Grado en Física y Matemáticas.
        \item[Grupo] Único.
        %\item[Profesor] José María Espinar García.
        \item[Descripción] Examen Ordinario 
        \item[Fecha] 23 de junio de 2020.
        % \item[Duración] 60 minutos.
    
    \end{description}
    \newpage

    \subsection*{TEORÍA (2.5 puntos)}
    \begin{ejercicio}
        Obtener la expresión analítica para la función de densidad del máximo y del mínimo de $n$ variables aleatorias independientes e identicamente distribuidas con función de distribución marginal $F$.\\

        Sean $X_1,\dots,X_n$ variables aleatorias independientes e idénticamente distribuidas con función de distribución $F_{X_i}=F$ para todo $i\in \{1,\dots,n\}$. Consideramos el vector aleatorio $X=(X_1,\dots,X_n)$.\\

        Respecto a la distribución del máximo, tenemos que:
        \begin{align*}
            F_{\max(X)}(x) &= P[X_1\leq x,\dots,X_n\leq x] = \prod_{i=1}^n P[X_i\leq x] = F(x)^n\qquad \forall x\in \mathbb{R}.
        \end{align*}

        Por tanto, su función de densidad es:
        \begin{equation*}
            f_{\max(X)}(x) = \frac{d}{dx}F_{\max(X)}(x) = nF(x)^{n-1}f(x)\qquad \forall x\in \mathbb{R}.
        \end{equation*}
        donde $f$ es la función de densidad común a todas las variables aleatorias, que además cumple $f=F'$.\\

        Respecto a la distribución del mínimo, tenemos que:
        \begin{align*}
            F_{\min(X)}(x) &= 1-P[X_1>x,\dots,X_n>x] = 1-\prod_{i=1}^n P[X_i>x] = 1-(1-F(x))^n\qquad \forall x\in \mathbb{R}.
        \end{align*}

        Por tanto, su función de densidad es:
        \begin{equation*}
            f_{\min(X)}(x) = \frac{d}{dx}F_{\min(X)}(x) = n(1-F(x))^{n-1}f(x)\qquad \forall x\in \mathbb{R}.
        \end{equation*}

    \end{ejercicio}

    \begin{ejercicio}
        Dada una sucesión $\{X_n\}_{n\in \mathbb{N}}$ de variables aleatorias independientes e idénticamente distribuidas tales que existen $E[X_n] = \mu$ y $E[X_n^2] < \infty$, probar que:
        $$\frac{S_n}{n}\stackrel{P}{\rightarrow} \mu$$

        Este resultado está probado como corolario directo de la Ley de Khintchine.
    \end{ejercicio}

    \subsection*{PROBLEMAS (7.5 puntos)}
    \begin{ejercicio}
        Una persona recibe un ticket que puede acumular para canjear por un premio si al elegir al azar dos números cualesquiera, $X$ entre 0 y 1 e $Y$ entre 1 y 3, su suma $X+Y<2$. Este procedimiento se puede repetir de forma indefinida e independiente. Calcular el número de elecciones necesarias para poder obtener un premio que requiere de al menos 10 tickets acumulados con probabilidad superior a $0.9495$.
        \begin{observacion}
            Aunque el enunciado no lo especifica, entendemos que $X$ e $Y$ son variables aleatorias continuas (podrían ser discretas).
        \end{observacion}

        Tenemos que:
        \begin{equation*}
            X\sim \cc{U}([0,1])\quad Y\sim \cc{U}([1,3])
        \end{equation*}

        Por tanto, sus funciones de densidad son:
        \begin{align*}
            f_X(x) &= \begin{cases}
                1 & \text{si } x\in [0,1] \\
                0 & \text{en otro caso}
            \end{cases} \\
            f_Y(y) &= \begin{cases}
                \nicefrac{1}{2} & \text{si } y\in [1,3] \\
                0 & \text{en otro caso}
            \end{cases}
        \end{align*}

        Buscamos calcular en primer lugar $P[X+Y<2]$. Para ello, por ser $X$ e $Y$ independientes, tenemos que:
        \begin{equation*}
            f_{X,Y}(x,y) = f_X(x)f_Y(y) = \begin{cases}
                \nicefrac{1}{2} & \text{si } (x,y)\in [0,1]\times [1,3] \\
                0 & \text{en otro caso}
            \end{cases}
        \end{equation*}

        Veamos el conjunto gráficamente, junto con la recta $x+y=2$:
        \begin{center}
            \begin{tikzpicture}
                \begin{axis}[
                    axis lines = center,
                    xlabel = $x$,
                    ylabel = $y$,
                    xmin = -1, xmax = 2,
                    ymin = -1, ymax = 4,
                    xtick = {0,1},
                    ytick = {0,1,2,3},
                    axis equal,
                ]
                    
                    \addplot[domain=-2:5, samples=2, color=red]{2-x};

                    % Rectangulo de x\in [0,1] e y\in [1,3]
                    \addplot[fill=orange, fill opacity=0.2] coordinates {(0,1) (1,1) (1,3) (0,3) (0,1)};

                    % Triangulo (0,1) (1,1) (0,2)
                    \addplot[fill=blue, fill opacity=0.2] coordinates {(0,1) (1,1) (0,2) (0,1)};
                \end{axis}
            \end{tikzpicture}
        \end{center}

        Por tanto, tenemos que:
        \begin{align*}
            P[X+Y<2] &= \int_{0}^{1}\int_{1}^{2-x}\frac{1}{2}dydx = \frac{1}{2}\int_{0}^{1}(2-x-1)dx = \frac{1}{2}\int_{0}^{1}(1-x)dx = \frac{1}{2}\left[x-\frac{x^2}{2}\right]_{0}^{1} =\\&= \frac{1}{2}\left(1-\frac{1}{2}\right) = \frac{1}{4}
        \end{align*}

        Por tanto, la probabilidad de obtener un ticket es $p=\nicefrac{1}{4}$. Sea ahora $Z_i$ la variable aleatoria que toma el valor 1 si se obtiene un ticket en la $i$-ésima elección y 0 en otro caso. Entonces, $Z_i\sim \cc{B}(1,p)$, con $p=\nicefrac{1}{4}$.\\

        Consideramos ahora $S_n=\sum\limits_{i=1}^n Z_i$, que por la propiedad reproductiva de la distribución binomial, sabemos que $S_n\sim \cc{B}(n,p)$. De hecho, $S_n$ es el número de tickets obtenidos en las $n$ elecciones. Por tanto, buscamos calcular $n$ tal que:
        \begin{equation*}
            P[S_n\geq 10] > 0.9495
        \end{equation*}

        Por ser $S_n\sim \cc{B}(n,p)$, tenemos que:
        \begin{equation*}
            E[S_n] = np\qquad \Var[S_n] = np(1-p)
        \end{equation*}

        Por tanto, buscamos $n$ tal que:
        \begin{align*}
            P[S_n\geq 10] > 0.9495
            &\Longleftrightarrow P\left[\dfrac{S_n-E[S_n]}{\sqrt{\Var[S_n]}}\geq \dfrac{10-np}{\sqrt{np(1-p)}}\right] > 0.9495 \\
            &\Longleftrightarrow P\left[\dfrac{S_n-E[S_n]}{\sqrt{\Var[S_n]}}\geq \dfrac{10-\nicefrac{n}{4}}{\sqrt{n}\cdot \nicefrac{\sqrt{3}}{4}}\right] > 0.9495
        \end{align*}

        Por el Teorema de Lévy, tenemos que:
        \begin{equation*}
            \left\{\dfrac{S_n-E[S_n]}{\sqrt{\Var[S_n]}}\right\}
            \stackrel{L}{\rightarrow} \cc{N}(0,1)
        \end{equation*}

        Por tanto, sea $Z\sim \cc{N}(0,1)$, buscamos $n$ tal que:
        \begin{equation*}
            P[Z\geq \dfrac{40-n}{\sqrt{n}\cdot \sqrt{3}}] > 0.9495
            \Longleftrightarrow
            P[Z\leq -\dfrac{40-n}{\sqrt{n}\cdot \sqrt{3}}] > 0.9495
        \end{equation*}

        Consultando la tabla de la distribución normal estándar, obtenemos que:
        \begin{align*}
            -\dfrac{40-n}{\sqrt{n}\cdot \sqrt{3}}>1.64&\Longleftrightarrow 40-n<-1.64\sqrt{n}\cdot \sqrt{3}
            \Longleftrightarrow n-2.84\sqrt{n}-40>0
        \end{align*}

        Resolvemos la igualdad:
        \begin{align*}
            \sqrt{n} &= \dfrac{2.84\pm \sqrt{2.84^2+4\cdot 40}}{2} = \dfrac{2.84\pm 12.96}{2} = 7.9\Longrightarrow n = 62.41
        \end{align*}

        Evaluando, vemos que como solución sirve cualquier $n>62.41$. Como se entiende que se pide el número entero más pequeño, la respuesta es $n=63$.\\

        Por tanto, debe realizar al menos 63 elecciones para obtener un premio con probabilidad superior a $0.9495$.
    \end{ejercicio}

    \begin{ejercicio}
        De un vector aleatorio bidimensional $(X,Y)$ se sabe que sus rectas de regresión son $5y-x+1=0$ y $2x-5y+2=0$:
        \begin{enumerate}
            \item Identificar la recta de regresión de $Y$ sobre $X$.
            
            Supongamos que la recta de regresión de $Y$ sobre $X$ es $y=\nicefrac{x}{5}-\nicefrac{1}{5}$, luego la recta de regresión de $X$ sobre $Y$ es $x=\nicefrac{5y}{2}-1$. Por tanto, como el coeficiente de determinación es el producto de las pendientes de las rectas de regresión, tenemos que:
            \begin{equation*}
                \rho_{X,Y}^2 = \frac{1}{5}\cdot \frac{5}{2} = \frac{1}{2}<1
            \end{equation*}
            Por tanto, el resultado es coherente y la recta de regresión de $Y$ sobre $X$ es:
            \begin{equation*}
                y=\frac{x}{5}-\frac{1}{5}
            \end{equation*}
            Si hubiésemos hecho la suposición contraria, habríamos obtenido un coeficiente de determinación mayor que 1, lo cual no tiene sentido.
            \item Obtener una medida de la proporción de varianza de cada variable que queda explicada por el modelo de regresión lineal.
            
            Esta medida es el coeficiente de determinación, que en este caso es $\rho_{X,Y}^2=\nicefrac{1}{2}$. Por tanto, el $50\%$ de la varianza de $Y$ queda explicada por el modelo de regresión lineal.
            \item Calcular la esperanza del vector bidimensional.
            
            Sabemos que este es el punto de corte de las rectas de regresión. Por tanto, resolviendo el sistema de ecuaciones:
            \begin{align*}
                5y-x+1 &= 0 \\
                2x-5y+2 &= 0
            \end{align*}
            obtenemos que $x=-3$ e $y=\nicefrac{-4}{5}$. Por tanto, tenemos que:
            \begin{equation*}
                E[(X,Y)] = (-3,-\nicefrac{4}{5})
            \end{equation*}
        \end{enumerate}
    \end{ejercicio}

    \begin{ejercicio}
        Dado el vector bidimensional $(X,Y)$ distribuido uniformemente en el recinto limitado $R = \{(x,y)\in \mathbb{R}^2 \mid -2\leq x \leq y \leq -x\}$:
        \begin{enumerate}
            \item Obtener su función de densidad conjunta.
            
            Representamos en primer lugar dicho conjunto:
            \begin{figure}[H]
                \centering
                \begin{tikzpicture}
                    \begin{axis}
                        [
                            axis lines = center,
                            xlabel = $X$,
                            ylabel = $Y$,
                            xmin = -3, xmax = 3,
                            ymin = -3, ymax = 3,
                        ]
                        \addplot[fill=red, opacity=0.3] coordinates {(-2, -2) (0, 0) (-2, 2)} -- cycle;
                        \node at (-1.5, 0.5) {$R$};

                        % R0: (-2,3) (-2,-2) (3,-2) (3,-3) (-3,-3) (-3,3) (-2,3)
                        \addplot[fill=blue, opacity=0.2] coordinates {
                            (-2, 3) (-2, -2) (3, -2) (3, -3) (-3, -3) (-3, 3) (-2, 3)
                        } -- cycle;
                        \node at (-2.5, -2.5) {$R_0$};

                        % R1: (-2,-2) (0,0) (3,0) (3,-2)
                        \addplot[fill=green, opacity=0.2] coordinates {
                            (-2, -2) (0, 0) (3, 0) (3, -2)
                        } -- cycle;
                        \node at (1.5, -1) {$R_1$};

                        % R2: (-2,2) (0,0) (0,2)
                        \addplot[fill=yellow, opacity=0.2] coordinates {
                            (-2, 2) (0, 0) (0, 2)
                        } -- cycle;
                        \node at (-0.8, 1.3) {$R_2$};

                        % R3: (0,2) (0,0) (3,0) (3,2)
                        \addplot[fill=orange, opacity=0.2] coordinates {
                            (0, 2) (0, 0) (3, 0) (3, 2)
                        } -- cycle;
                        \node at (1.5, 1) {$R_3$};

                        % R4: (-2,2) (0,2) (0,3) (-2,3)
                        \addplot[fill=purple, opacity=0.2] coordinates {
                            (-2, 2) (0, 2) (0, 3) (-2, 3)
                        } -- cycle;
                        \node at (-1, 2.5) {$R_4$};

                        % R5: (0,2) (3,2) (3,3) (0,3)
                        \addplot[fill=gray, opacity=0.2] coordinates {
                            (0, 2) (3, 2) (3, 3) (0, 3)
                        } -- cycle;
                        \node at (1.5, 2.5) {$R_5$};
                    \end{axis}
                \end{tikzpicture}
            \end{figure}

            Tenemos que, para $x\in [-2, 0]$ y $y\in [x, -x]$, la función de densidad conjunta es:
            \begin{equation*}
                f_{(X,Y)}(x, y) = k,\qquad k\in \bb{R}^+
            \end{equation*}

            Para que sea una función de densidad, hemos de tener que:
            \begin{equation*}
                1 = \int_{\bb{R}^2} f = \int_R f = \int_R k
            \end{equation*}

            Tenemos dos opciones:
            \begin{description}
                \item[Integración Normal:]
                \begin{align*}
                    1 = \int_R k &= \int_{-2}^0 \int_{x}^{-x} k \, dy \, dx = \int_{-2}^0 k(-2x) \, dx =\\&= -2k\left[\dfrac{x^2}{2}\right]_{-2}^0 = -2k\left[0-2\right] = 4k \Longrightarrow k = \dfrac{1}{4}
                \end{align*}

                \item[Razonando según la Forma de $R$:]
                \begin{align*}
                    1 = \int_R k = k\lm(R) = k\cdot \dfrac{4\cdot 2}{2} = 4k \Longrightarrow k = \dfrac{1}{4}
                \end{align*}
            \end{description}
            En cualquier caso, para el valor de $k=\nicefrac{1}{4}$, la función de densidad conjunta es integrable, no negativa e integra $1$.
            \item Obtener su función de distribución conjunta.
            Distinguimos casos:
            \begin{itemize}
                \item \ul{Si $x < -2$ o $y < -2$} (Zona $R_0$):
                \begin{equation*}
                    F_{(X,Y)}(x, y) \int_{-\infty}^x \int_{-\infty}^y f_{(X,Y)}(u, v) \, dv \, du = 0
                \end{equation*}

                \item \ul{Si $x\in [-2, 0]$ y $y\in [x, -x]$} (Zona $R$):
                \begin{align*}
                    F_{(X,Y)}(x, y) &= \int_{-\infty}^x \int_{-\infty}^y f_{(X,Y)}(u, v) \, dv \, du = \int_{-2}^x \int_{u}^{y} \dfrac{1}{4} \, dv \, du
                    =\\&= \int_{-2}^x \dfrac{1}{4}(y-u) \, du = \dfrac{1}{4}\left[yu - \dfrac{u^2}{2}\right]_{-2}^x = \dfrac{1}{4}\left[yx - \dfrac{x^2}{2} + 2y+2\right]
                \end{align*}

                \item \ul{Si $y<0$ y $x > y$} (Zona $R_1$):
                \begin{align*}
                    F_{(X,Y)}(x, y) &= \int_{-\infty}^x \int_{-\infty}^y f_{(X,Y)}(u, v) \, dv \, du = \int_{-2}^y \int_{u}^{y} \dfrac{1}{4} \, dv \, du
                    =\\&= \int_{-2}^y \dfrac{1}{4}(y-u) \, du = \dfrac{1}{4}\left[yu - \dfrac{u^2}{2}\right]_{-2}^y = \dfrac{1}{4}\left[y^2 - \dfrac{y^2}{2} + 2y+2\right]
                    =\\&= \dfrac{1}{4}\left[\dfrac{y^2}{2} + 2y + 2\right]
                \end{align*}

                \item \ul{Si $x\in [-2, 0]$ y $y\in [-x, 2]$} (Zona $R_2$):
                \begin{align*}
                    F_{(X,Y)}(x, y) &= \int_{-\infty}^x \int_{-\infty}^y f_{(X,Y)}(u, v) \, dv \, du = \int_{-2}^{-y} \int_{u}^{y} \dfrac{1}{4} \, dv \, du+ \int_{-y}^{x} \int_{u}^{-u} \dfrac{1}{4} \, dv \, du
                    =\\&= \int_{-2}^{-y} \dfrac{1}{4}(y-u) \, du + \int_{-y}^{x} \dfrac{1}{4}(-2u) \, du = \dfrac{1}{4}\left[yu - \dfrac{u^2}{2}\right]_{-2}^{-y} + \dfrac{1}{4}\left[-u^2\right]_{-y}^{x}
                    =\\&= \dfrac{1}{4}\left[\cancel{y(-y)} - \dfrac{(-y)^2}{2} + 2y+2\right] + \dfrac{1}{4}\left[-x^2 + \cancel{y^2}\right]
                    =\\&= \dfrac{1}{4}\left[2y + 2 - \dfrac{y^2}{2} - x^2\right]
                \end{align*}

                \item \ul{Si $y\in [0, 2]$ y $x \geq 0$} (Zona $R_3$):
                \begin{align*}
                    F_{(X,Y)}(x, y) &= \int_{-\infty}^x \int_{-\infty}^y f_{(X,Y)}(u, v) \, dv \, du = \int_{-2}^{-y} \int_{u}^{y} \dfrac{1}{4} \, dv \, du+ \int_{-y}^{0} \int_{u}^{-u} \dfrac{1}{4} \, dv \, du
                    =\\&= \int_{-2}^{-y} \dfrac{1}{4}(y-u) \, du + \int_{-y}^{0} \dfrac{1}{4}(-2u) \, du = \dfrac{1}{4}\left[yu - \dfrac{u^2}{2}\right]_{-2}^{-y} + \dfrac{1}{4}\left[-u^2\right]_{-y}^{0}
                    =\\&= \dfrac{1}{4}\left[\cancel{y(-y)} - \dfrac{(-y)^2}{2} + 2y+2\right] + \dfrac{1}{4}\left[0 - \cancel{y^2}\right]
                    =\\&= \dfrac{1}{4}\left[2y + 2 - \dfrac{y^2}{2}\right]
                \end{align*}

                \item \ul{Si $x\in [-2, 0]$ y $y\geq 2$} (Zona $R_4$):
                \begin{align*}
                    F_{(X,Y)}(x, y) &= \int_{-\infty}^x \int_{-\infty}^y f_{(X,Y)}(u, v) \, dv \, du = \int_{-2}^{x} \int_{u}^{-u} \dfrac{1}{4} \, dv \, du
                    =\\&= \int_{-2}^{x} \dfrac{1}{4}(-2u) \, du = \dfrac{1}{4}\left[-u^2\right]_{-2}^{x} = \dfrac{1}{4}\left[-x^2 + 4\right]
                \end{align*}

                \item \ul{Si $x\geq 0, y\geq 2$} (Zona $R_5$):
                \begin{equation*}
                    F_{(X,Y)}(x, y) = 1
                \end{equation*}
            \end{itemize}

            Por tanto, la función de distribución conjunta es:
            \begin{equation*}
                F_{(X,Y)}(x, y) = \begin{cases}
                    0 & \text{si } x < -2 \text{ o } y < -2\\
                    \dfrac{1}{4}\left[yx - \dfrac{x^2}{2} + 2y+2\right] & \text{si } x\in [-2, 0] \text{ y } y\in [x, -x]\\
                    \dfrac{1}{4}\left[\dfrac{y^2}{2} + 2y + 2\right] & \text{si } y<0 \text{ y } x > y\\
                    \dfrac{1}{4}\left[2y + 2 - \dfrac{y^2}{2} - x^2\right] & \text{si } x\in [-2, 0] \text{ y } y\in [-x, 2]\\
                    \dfrac{1}{4}\left[2y + 2 - \dfrac{y^2}{2}\right] & \text{si } y\in [0, 2] \text{ y } x \geq 0\\
                    \dfrac{1}{4}\left[-x^2 + 4\right] & \text{si } x\in [-2, 0] \text{ y } y\geq 2\\
                    1 & \text{si } x\geq 0, y\geq 2
                \end{cases}
            \end{equation*}
            \item Obtener la probabilidad de que $X+Y+1\geq 0$.
            Veamos qué conjunto representa $X + Y + 1 \geq 0$:
            \begin{figure}[H]
                \centering
                \begin{tikzpicture}
                    \begin{axis}
                        [
                            axis lines = center,
                            xlabel = $X$,
                            ylabel = $Y$,
                            xmin = -3, xmax = 3,
                            ymin = -3, ymax = 3,
                        ]
                        \addplot[name path=A, fill=red, opacity=0.3] coordinates {(-2, -2) (0, 0) (-2, 2)} -- cycle;

                        % Recta y = -x-1
                        \addplot[name path=B, domain=-3:2.5, samples=2, dashed] {-x-1};

                        % Add a node to label the line
                        \node at (axis cs:1,-2) [anchor=west] {$y=-x-1$};

                        % R0: (-2,2) (0,0) (-0.5,-0.5) (-2, 1.5)
                        \addplot[fill=blue, opacity=0.2] coordinates {
                            (-2, 2) (0, 0) (-0.5, -0.5) (-2, 1)
                        } -- cycle;
                    \end{axis}
                \end{tikzpicture}
            \end{figure}

            Tenemos por tanto que:
            \begin{align*}
                P[X + Y + 1 \geq 0] &= \int_{-2}^{\nicefrac{-1}{2}} \int_{-x-1}^{-x} \dfrac{1}{4} \, dy \, dx + \int_{\nicefrac{-1}{2}}^{0} \int_{x}^{-x} \dfrac{1}{4} \, dy \, dx
                =\\&= \int_{-2}^{\nicefrac{-1}{2}} \dfrac{1}{4}(-x+1+x) \, dx + \int_{\nicefrac{-1}{2}}^{0} \dfrac{1}{4}(-2x) \, dx
                =\\&= \int_{-2}^{\nicefrac{-1}{2}} \dfrac{1}{4} \, dx + \int_{\nicefrac{-1}{2}}^{0} -\dfrac{1}{2}x \, dx
                =\\&= \dfrac{1}{4}\left[x\right]_{-2}^{\nicefrac{-1}{2}} - \dfrac{1}{2}\left[\dfrac{x^2}{2}\right]_{\nicefrac{-1}{2}}^{0}
                = \dfrac{1}{4}\left[\dfrac{-1}{2}+2\right] - \dfrac{1}{2}\left[0 - \dfrac{1}{8}\right] = \dfrac{7}{16}
            \end{align*}
        \end{enumerate}
    \end{ejercicio}

    \begin{ejercicio}
        Dado el vector bidimensional $(X,Y)$ con la siguiente función masa de probabilidad conjunta:
        \begin{table}[H]
        \centering
        \begin{tabular}{|c|ccc|}
            \hline
            $X\backslash Y$ & 0 & 1 & 2 \\
            \hline
            1 & $\nicefrac{1}{4}$ & 0 & 0 \\
            2 & 0 & $\nicefrac{1}{4}$ & 0 \\
            3 & $\nicefrac{1}{4}$ & 0 & $\nicefrac{1}{4}$ \\
            \hline
        \end{tabular}
        \end{table}
        \begin{enumerate}
            \item Obtener la mejor aproximación minimo-cuadrática a la variable $Y$ conocidos valores de la variable $X$, así como calcular una medida de la bondad del ajuste.
            
            Completamos en primer lugar la función masa de probabilidad conjunta calculando las marginales:
            \begin{table}[H]
                \centering
                \begin{tabular}{c|ccc||c}
                    $X\backslash Y$ & 0 & 1 & 2 \\
                    \hline
                    1 & $\nicefrac{1}{4}$ & 0 & 0 & $\nicefrac{1}{4}$\\
                    2 & 0 & $\nicefrac{1}{4}$ & 0 & $\nicefrac{1}{4}$ \\
                    3 & $\nicefrac{1}{4}$ & 0 & $\nicefrac{1}{4}$ & $\nicefrac{1}{2}$ \\ \hline
                    \hline
                    & $\nicefrac{1}{2}$ & $\nicefrac{1}{4}$ & $\nicefrac{1}{4}$
                \end{tabular}
                \end{table}
            
            Mostramos ahora en la siguiente tabla la distribución condicionada a $X$, es decir, $P[Y|X=x]$:
            \begin{table}[H]
            \centering
            \begin{tabular}{c|ccc}
                $X\backslash Y$ & 0 & 1 & 2 \\
                \hline
                1 & $1$ & 0 & 0 \\
                2 & 0 & $1$ & 0 \\
                3 & $\nicefrac{1}{2}$ & 0 & $\nicefrac{1}{2}$
            \end{tabular}
            \end{table}

            Por tanto, tenemos que:
            \begin{align*}
                E[Y|X=1] &= 0\cdot 1 + 1\cdot 0 + 2\cdot 0 = 0 \\
                E[Y|X=2] &= 0\cdot 0 + 1\cdot 1 + 2\cdot 0 = 1 \\
                E[Y|X=3] &= 0\cdot \nicefrac{1}{2} + 1\cdot 0 + 2\cdot \nicefrac{1}{2} = 1
            \end{align*}

            Por tanto, la mejor aproximación minimo-cuadrática a la variable $Y$ conocidos valores de la variable $X$ es:
            \begin{equation*}
                \hat{Y} = \begin{cases}
                    0 & \text{si } X=1 \\
                    1 & \text{si } X=2,3
                \end{cases}
            \end{equation*}

            Calculamos ahora una medida de la bondad del ajuste. Calculamos previamente:
            \begin{align*}
                E[Y]&=0\cdot \nicefrac{1}{2} + 1\cdot \nicefrac{1}{4} + 2\cdot \nicefrac{1}{4} = \nicefrac{3}{4}\\
                E[Y^2]&=0^2\cdot \nicefrac{1}{2} + 1^2\cdot \nicefrac{1}{4} + 2^2\cdot \nicefrac{1}{4} = \nicefrac{5}{4}\\
                \Var[Y] &= E[Y^2] - E[Y]^2 = \nicefrac{5}{4} - \nicefrac{9}{16} = \nicefrac{11}{16}\\
                E[(E[Y|X])^2] &= \sum_{x\in \{1,2,3\}} E[Y|X=x]^2P[X=x] = 0^2\cdot \nicefrac{1}{4} + 1^2\cdot \nicefrac{1}{4} + 1^2\cdot \nicefrac{1}{2} = \nicefrac{3}{4}\\
                E[E[Y|X]] &= E[Y] = \nicefrac{3}{4}\\
                \Var[E[Y|X]] &= E[(E[Y|X])^2] - E[E[Y|X]]^2 = \nicefrac{3}{4} - \nicefrac{9}{16} = \nicefrac{3}{16}
            \end{align*}

            Por tanto, la medida de la bondad del ajuste es:
            \begin{align*}
                \eta_{Y/X}^2 &= \dfrac{\Var[E[Y|X]]}{\Var[Y]} = \dfrac{\nicefrac{3}{16}}{\nicefrac{11}{16}} = \dfrac{3}{11}
            \end{align*}

            \item Obtener la aproximación lineal minimo-cuadrática de la variable $Y$ para $X=~1$.
            
            Podríamos pensar que se trata de la mejor aproximación minimo-cuadrática de la variable $Y$ para $X=1$, que sería $E[Y\mid X=1] = 0$. Sin embargo, esta no es la aproximación lineal. La aproximación lineal minimo-cuadrática de la variable $Y$ sobre $X$ es:
            \begin{equation*}
                Y-E[Y] = \dfrac{\Cov[X,Y]}{\Var[X]}(X-E[X])
            \end{equation*}

            Calculamos por tanto dichos valores:
            \begin{align*}
                E[X]&=1\cdot \nicefrac{1}{4} + 2\cdot \nicefrac{1}{4} + 3\cdot \nicefrac{1}{2} = \nicefrac{9}{4} \\
                E[X^2]&=1^2\cdot \nicefrac{1}{4} + 2^2\cdot \nicefrac{1}{4} + 3^2\cdot \nicefrac{1}{2} = \nicefrac{23}{4} \\
                \Var[X] &= E[X^2] - E[X]^2 = \nicefrac{23}{4} - \nicefrac{81}{16} = \nicefrac{11}{16} \\
                E[XY] &= 2\cdot 1\cdot \nicefrac{1}{4}+3\cdot 2\cdot \nicefrac{1}{4}=2\\
                \Cov[X,Y] &= E[XY] - E[X]E[Y] = 2 - \nicefrac{9}{4}\cdot \nicefrac{3}{4} = \nicefrac{5}{16}
            \end{align*}

            Por tanto, la aproximación lineal minimo-cuadrática de la variable $Y$ es:
            \begin{equation*}
                Y-\frac{3}{4} = \dfrac{\nicefrac{5}{16}}{\nicefrac{11}{16}}\left(X-\frac{9}{4}\right)
                \Longleftrightarrow
                Y = \frac{5}{11}X- \frac{3}{11}
            \end{equation*}

            Por tanto, evaluando en $X=1$, obtenemos que la aproximación lineal minimo-cuadrática de la variable $Y$ para $X=1$ es $$Y=\frac{2}{11}$$
            \item Obtener el error cuadrático medio para la aproximación del primer apartado.
            
            Tenemos que:
            \begin{align*}
                \text{E.C.M.}(E[Y\mid X])
                = E[\Var[Y\mid X]] = \Var[Y]-\Var[E[Y\mid X]] = \frac{11}{16} - \frac{3}{16} = \frac{8}{16} = \frac{1}{2}
            \end{align*}
        \end{enumerate}
    \end{ejercicio}

\end{document}
